\subsection{network namespace separato per la
verifica}\label{network-namespace-separato-per-la-verifica}

In blueagent, abbiamo introdotto OVN per semplificare la gestione di due
reti con stesso gateway.

Per evitare di aggiungere la complessita computazionale di OVN, si
puo\textquotesingle{} pensare di creare network namespace separato al
momento della verifica. All\textquotesingle interno di questo nuovo
namespace, chiamato \(\text{verify-}k\) si crea la rete di verifica
composta dal solo bridge L2 e si assegna a questo bridge il gateway
originale.

I container duplicati dovranno poi essere collegati al nuovo bridge L2.

Per consentire l\textquotesingle accesso alla rete bisogna usare una
coppia di interfacce \passthrough{\lstinline!veth!}, per collegare il
namespace nuovo al root namespace. L\textquotesingle inoltro dei
pacchetti in questo contesto avviene applicando regole di forwarding e
usando il masquerading.

In questo modo i due bridge con stesso ip gateway possono coesistere
poiche riesiedono in namespace differenti. In questo contesto ambiguo ci
pensa il kernel ad effettuare il routing dei pacchetti verso le giuste
interfacce.

Tutta via l\textquotesingle implementazione non e\textquotesingle{}
immediata ed l\textquotesingle uso di Incus non si sposa bene con la
creazione di nuovi network namespace dato che Incus non fornisce
un\textquotesingle api apposita. Dunque bisogna gestire manualmente le
interfacce di rete raggirando Incus.

Quando eseguo il container usando l\textquotesingle interfaccia di
Incus, non posso collegare le interfacce di rete dei container al bridge
di verifica dentro il namespace \(\text{verify-}k\). Questo
principalmente perche\textquotesingle{} il motore di Incus esegue su un
namespace separato, dunque e\textquotesingle{} impossibile che lo veda.

Bisogna dunque raggirare incus, ed istanziare coppie di \emph{veth} che
gestiscano questo disallineamento. Ma raggirare incus vuol dire che
incus non ha idea dell\textquotesingle esistenza di queste interfacce,
creando un disallinamento tra container ed interfacce di rete, che nello
sviluppo puo\textquotesingle{} diventare problematico da gestire.

Dal lato del codice, aumenta la complessita: invece di fare richieste
verso l\textquotesingle API di incus per gestire le interfacce, devo
eseguire vari comandi per gestire il routing e le coppie \emph{veth}.

\begin{center}\rule{0.5\linewidth}{0.5pt}\end{center}

\subsubsection{Conclusione}\label{conclusione}

L\textquotesingle uso di network namespace dedicati
sull\textquotesingle host sostituisce l\textquotesingle adozione di un
software SDN strutturato con una complessa selva di script manuali di
routing e manipolazione delle veth.

L\textquotesingle approccio è stato quindi scartato:
l\textquotesingle integrazione nativa di OVN con Incus offre un
isolamento di rete standardizzato, manutenibile via API e già
predisposto per ambienti distribuiti su cluster.
